\documentclass[10pt]{article}
\usepackage[margin=0.85in]{geometry}
\usepackage{amsmath,amssymb}
\usepackage[hidelinks]{hyperref}

\begin{document}
\begin{center}
  {\Large\bfseries A non-tree graph with a formal configuration space\par}
  \smallskip
  Disproof of Conjecture 00000004113
\end{center}

\noindent
The cycle $C_3$ is a non-tree graph whose two-particle configuration
space is rationally formal. The unordered space deformation retracts
onto a circle. The ordered space also has the homotopy type of a circle.

\section*{The unordered configuration space}

The graph $C_3$ has vertices $0,1,2$ and edges $01,12,20$. Let $U_2(C_3)$
be the space of unordered pairs of distinct points of its geometric
realization. Its discretized subspace $UD_2(C_3)$ uses pairs of open graph
cells with disjoint closures, with the two factors identified under
interchange.

The subdivision criterion in Farley--Sabalka~\cite[Theorem~3.1]{fs},
proved in the strengthened form by Prue--Scrimshaw~\cite[Theorem~3.2]{ps},
applies at $n=2$. There are at least two vertices. The paths connecting
distinct vertices of degree other than two would need at least one edge,
and every vertex here has degree two. A homotopically nontrivial closed
path has length at least three. Indeed, there are no loops and any
closed walk of length two backtracks. Hence $U_2(C_3)$ deformation
retracts onto $UD_2(C_3)$.

The zero-cells of $UD_2(C_3)$ are the three unordered pairs of distinct
original vertices. Identify each pair with its missing vertex. A one-cell
consists of an original edge $ab$ and the remaining vertex $c$. Its
endpoints are $\{c,a\}$ and $\{c,b\}$, which identify with $b$ and $a$.
Every original edge gives exactly one such cell. There are no two-cells
because any two original edges have a common endpoint. These vertex and
edge identifications give a graph isomorphism $UD_2(C_3)\cong C_3$, whose
geometric realization is $S^1$.

\section*{Formality and the ordered space}

Take a closed element $\omega\in A_{\mathrm{PL}}^1(S^1)$ representing a
generator of $H^1(S^1;\mathbb Q)$. Sullivan's rational polynomial forms
compute rational cohomology~\cite{fh}. Graded commutativity gives
$\omega^2=0$. The map $\Lambda_{\mathbb Q}(a)\to A_{\mathrm{PL}}(S^1)$
sending $a$ to $\omega$, where $|a|=1$ and $da=0$, induces an isomorphism
in every cohomological degree. Thus $S^1$ is formal, as is $U_2(C_3)$ by
homotopy invariance.

For the ordered version, identify $|C_3|$ with the unit circle in
$\mathbb C$. The map $(z,w)\mapsto (z,z^{-1}w)$ is a homeomorphism from
the ordered configuration space to $S^1\times(S^1\setminus\{1\})$, with
inverse $(z,u)\mapsto(z,zu)$. The second factor is contractible, so the
ordered space is also formal. Either convention gives a counterexample to
the implication that formality forces the graph to be a tree.

\section*{Lean verification}

The Lean 4 project checks the original three-cycle and computes the
admissible cell pairs from their closures. It proves that there are no
square cells and constructs inverse vertex and edge maps preserving the
unordered incidences. The theorem
\texttt{C4113.cycle\_configuration\_certificate} records these finite
claims. The deformation retraction and rational model above give the
topological conclusion. Run \texttt{lake build} in the \texttt{lean}
directory to check the certificate.

\begin{thebibliography}{3}
\bibitem{fs}
D.~Farley and L.~Sabalka, \emph{Discrete Morse theory and graph braid groups},
Algebraic \& Geometric Topology \textbf{5} (2005), 1075--1109.
\href{https://doi.org/10.2140/agt.2005.5.1075}{DOI 10.2140/agt.2005.5.1075}.

\bibitem{ps}
P.~Prue and T.~Scrimshaw, \emph{Abrams's stable equivalence for graph braid groups},
\href{https://arxiv.org/abs/0909.5511}{arXiv:0909.5511}.

\bibitem{fh}
Y.~F\'elix and S.~Halperin, \emph{Rational homotopy via Sullivan models and enriched Lie algebras},
EMS Surveys in Mathematical Sciences \textbf{10} (2023), 101--122.
\href{https://doi.org/10.4171/EMSS/67}{DOI 10.4171/EMSS/67}.
\end{thebibliography}

\end{document}
